\cotefile{156-9}
\scanpage{58}
\mspage{54}
% archive sheet number: 57

satisfaisant les deux seules (?) conditions

\uline{$\unsure{\mathrm{Atspat}|\mathrm{At}_{\mathrm{spat}}}\ 1'$} Si $X\in A$, et $Y\trianglelefteq X$, alors $Y\in A$
($A$ fermé dans $\mathrm{Atspat}(\Sigma)$, pour $\trianglelefteq$).

\deleted{\uline{$\mathrm{Atspat}\ 2$}\quad $\forall S\in\Sigma$, \unsure{soit}}
\deleted{$\unsure{A}$} Soit
\[
\deleted{A_S=\{X\in A\mid\unsure{\widetilde X^{\unsure{\mathrm{o}|\circ|0}}|X^{\mathrm{o}}}\leq S\}\unsure{\Longleftrightarrow|\longleftarrow}\forall Y\in\widetilde X,\ \unsure{\text{où}}\ Y\leq S}
\]
% A redirection stroke leads down to the following cancelled formula.
\[
\deleted{\Sigma_0=\{S\in\Sigma\mid\exists X\in A\ \unsure{\vdash|+}}\quad X^{\mathrm{o}}\unsure{\mid|/}X\,.
\]

\uline{$\mathrm{Atspat}\ 2'$}\quad \deleted{$\unsure{\forall|S}$} Soit
\[
\Sigma'=\{X^{\mathrm{o}}\mid X\in A\}\subset\Sigma\,.
\]
% The next two lines are grouped by a vertical stroke in the left margin.
Pour \deleted{$\unsure{\forall}$} $S\in\Sigma$, soit
$\Sigma'_{\leq S}=\{\unsure{\zeta|\xi}\in\Sigma'\mid\zeta\leq S\}$.
Ceci posé, on a, pour $\forall S\in\Sigma$
\[
S\unsure{\equiv|=}\operatorname{Sup}\Sigma'_{\leq S}\,.
\]

Je m'attends que, pour de telles données, \deleted{l'atelier $A$} \unsure{partie} en question $A$ de $\mathrm{Atspat}(\Sigma)$, munie de la structure d'atelier induite, est un atelier spatial, et que l'on a un iso. can.
\[
\Sigma\simeq\Sigma_A
\]
moyennant l'application canonique
\[
A\xrightarrow{\unsure{\mathrm{D\acute ep}|\mathrm{Dup}}}\mathrm{Atspat}(\Sigma_A)
\]
devient l'\uline{inclusion} de $A$ dans $\mathrm{Atspat}(\Sigma)$.

\scanpage{59}
\mspage{55}
% archive sheet number: 58

9 \uline{Juillet}

Il sera plus commode de vérifier \inserted{d'abord}, en termes des données $\trianglelefteq_A$, $\unsure{|\mathrm{o}|_A|\lvert\circ\rvert_A}$ induites par $\unsure{\mathrm{Fig\acute elspat}}(\Sigma)$, les axiomes $\unsure{\mathrm{Atspat}|\mathrm{At}_{\mathrm{spat}}}\ 1-3$.
$\mathrm{Atspat}\ 1,2$ triviales. Pour $\mathrm{Atspat}\ 3$, on va déterminer \inserted{le \unsure{spatio}} $\Sigma_A$ (défini en termes de $|\mathrm{o}|_A$ seulement), et l'application
$A\overset{\unsure{\mathrm{D\acute ep}|\mathrm{Dup}}}{\mathrel{\text{---}}}\mathrm{Fig}\deleted{\unsure{\mathrm{\acute el}}}\unsure{\mathrm{spat\acute el}|\mathrm{spatial}}(\Sigma_A)$
(définie en termes de $|\mathrm{o}|_A$ \uline{et} $\trianglelefteq_A$), et voir que ça s'identifie resp. à $\Sigma$ et à l'inclusion tautologique $A\subset\mathrm{Figspat\acute el}(\unsure{\Sigma_A|\Sigma})$.

\deleted{\unsure{Soit}} Considérons les applications
\diagram{
\begin{tikzcd}[ampersand replacement=\&,column sep=large,row sep=small]
A \arrow[r,"\alpha"] \& \Sigma \arrow[r,"\beta"] \& \unsure{\mathcal P|\mathscr P|P}(A) \\
\& \unsure{S} \arrow[r,mapsto] \& \{X\in A\mid X^{\unsure{\mathrm{o}|\circ|0}}\leq S\} \\
X=(\Phi,\rho) \arrow[r,mapsto] \& \begin{gathered}X^{\mathrm{o}}=\text{plus}\\\text{\unsure{grd} él. de}\\\Phi\text{ pour }\rho\end{gathered} \&
\end{tikzcd}}
On a
\[
\beta(\alpha(X))=\{Z\in A\mid Z^{\mathrm{o}}\leq X^{\mathrm{o}}\}\subset\operatorname{supp}^{\mathrm{o}}_A X
\]
\deleted{$\unsure{\operatorname{supp}^{\mathrm{o}}(X)}$} et
\[
\overline\alpha\beta(S)=\{X^{\mathrm{o}}\mid X\in A,\ \text{t.q. }X^{\mathrm{o}}\leq S\}\subset\operatorname{supp}^{\mathrm{o}}_\Sigma(S)
\]
% slanted
\margin{Cf \unsure{sorite|truc} \unsure{explicité} plus bas, sur la \unsure{fonctorialité} de $\unsure{\mathbb\Sigma|\Sigma}$\ldots}
d'où résulte que $\alpha$ et $\beta$ induisent des bijections inverses l'une de l'autre entre $\mathbb\Sigma_A$ et $\mathbb\Sigma_\Sigma$. Utilisant
\diagram{
\begin{tikzcd}[ampersand replacement=\&,column sep=large]
\mathbb\Sigma_\alpha:\quad\mathbb\Sigma_A \arrow[r] \& \mathbb\Sigma_\Sigma \& \Sigma \arrow[l,"\operatorname{supp}^{\mathrm{o}}_\Sigma"']
\end{tikzcd}}
\margin{cf plus haut p.\ 48}
,
\deleted{on trouve l'application}
pour identifier $\mathbb\Sigma_A$ et $\Sigma$, et interpréter

\scanpage{60}
\mspage{56}
% archive sheet number: 59

en conséquence $\unsure{\mathrm{D\acute ep}|\mathrm{Dep}}_A$ comme
% slanted
\margin{(*) pour l'instant, on prend $\mathrm{D\acute ep}_A X$ \uline{sans} sa \unsure{relation} d'ordre, soit $\mathrm{D\acute ep}^{!}_A X$}
\diagram{
\begin{tikzcd}[ampersand replacement=\&,column sep=large,row sep=small]
A \arrow[r,"\mathrm{D\acute ep}^{!}_A"] \&
\begin{gathered}\inserted{\unsure{\mathcal P|\mathscr P|P}(\unsure{\mathbb\Sigma|\Sigma}_A)}^{(*)}\\
\deleted{\unsure{\mathcal P\mathcal P}}\underset{\rule{0.4pt}{3pt}\rule{2.4em}{0.4pt}\rule{0.4pt}{3pt}}{\deleted{(\unsure{\mathbb\Sigma_A})}}
\end{gathered}\arrow[r] \& \mathcal P\deleted{\unsure{\mathcal P}}(\Sigma)\\
\begin{gathered}X\\\Vert\\(\Phi,\rho)\end{gathered}\arrow[r,mapsto] \&
\{\operatorname{supp}^{\unsure{\mathrm{o}|\circ|0}}_A(X')\mid X'\in\underset{\substack{\Vert\\\Phi}}{\widetilde X}\}\arrow[r,"?"] \& \deleted{\unsure{\mathrm{cf}|\Phi}}
\end{tikzcd}}
Or $\operatorname{supp}^{\mathrm{o}}_A(X')$\inserted{$\in\deleted{\unsure{\mathcal P}}\mathbb\Sigma_A\unsure{)}$} devient, pour $X'\in\Phi$\inserted{$\subset\Sigma$} donné, par $\mathbb\Sigma_A\xrightarrow{\sim}\mathbb\Sigma_\Sigma$ $\operatorname{supp}^{\mathrm{o}}_\Sigma(X')$, et par $\mathbb\Sigma_\Sigma\leftarrow\Sigma$, $X'$.
Donc l'image de $X=(\Phi,\rho)$ par \deleted{$\unsure{A}$} le composé $(*)$ n'est autre que $\Phi$, et la relation d'ordre à mettre dessus est celle déduite par transport de structure de celle $\rho$ de $\widetilde X=\Phi$, par l'application identique de $\widetilde X\xrightarrow{\sim}\Phi$, i.e. c'est $\rho$.
On a gagné!

\uline{Scholie} La donnée d'un atelier spatial revient à celle d'un spatios $(\Sigma,\leq\inserted{\text{ou }\overset{\mathrm{o}}{\ll}},\unsure{C|\mathcal C})$, muni de la structure supplémentaire définie par un $A\subset\unsure{\mathrm{Figspat\acute el}|\mathrm{Figspatial}}(\Sigma)$, satisfaisant les deux conditions $\unsure{\mathrm{Atspat}|\mathrm{At}_{\mathrm{spat}}}\ 1',2'$ p.\ 54. Ainsi l'ensemble de base de la

\scanpage{61}
\mspage{57}
% archive sheet number: 60

structure envisagée (qui avait été, tour à tour, $\unsure{\mathscr F_\xi|\mathcal F_\xi|F_\xi|\mathscr F_\zeta}$ (mais en \unsure{sens} inverse!), $A$ ou $(A,\mathscr F)$, $\unsure{\mathscr L|\mathcal L|L}$, est à présent $\Sigma$, le ``spatios des supports''. [Ce dernier est particulièrement commode pour formuler les axiomes des supports, cf plus bas, même pour des ateliers pas \unsure{néc\textsuperscript{t}} spatiaux.]

\uline{Sorite} du \deleted{$\unsure{\mathbb\Sigma|\Sigma}$} spatios $\mathbb\Sigma_{\deleted{\unsure{X|E}}\underline E}\unsure{)}$ associé à un ens. $\unsure{E|\mathcal E}$ muni d'\inserted{\unsure{une}} une relation symétrique anti\inserted{\unsure{ré}}flexive (notée $\unsure{|\mathrm{o}|_X|\lvert\circ\rvert_X}$), \unsure{i.e. qu'il s'agit} d'une relation telle que
\[
x|\mathrm{o}|x\Longrightarrow\forall y\in E,\quad x|\mathrm{o}|y\,.
\]
% slanted
% The words read as plein-vide are above/right of catégorie; their attachment is uncertain.
\margin{\uline{NB} \unsure{La catégorie} \unsure{plein-vide} \unsure{des ens. disjoints} \deleted{\unsure{pointés}} \unsure{est} \uline{\unsure{isomorphe}} à celle des \unsure{graphes stricts}, en posant $x|\mathrm{o}|x$ \unsure{ssi} $\forall y(\ne x),\ \{x,y\}\unsure{\notin|\in}L$.}

\unsure{1°)|0°)} Soient $\underline E=(E,|\mathrm{o}|_E)$, $\underline E'=(E',|\mathrm{o}|_{E'})$ munis de relations symétriques, une correspondance
\[
E\xrightarrow{\unsure{\varphi|\phi}}\unsure{\mathcal P|\mathscr P|P}(E')
\]
est dite un \uline{morphisme} de $\underline E$ dans $\underline E'$\inserted{(\unsure{i.e.} \uline{compatible} \unsure{à} $|\mathrm{o}|_E,|\mathrm{o}|_{E'}$)}, si
\[
\deleted{\forall\unsure{x},y\in E,}
\]
\[
\forall x,y\in E,\ \text{et }x'\in\varphi(x),\ y'\in\varphi(y),\quad x|\mathrm{o}|y\Longrightarrow x'|\mathrm{o}|y'
\]
% slanted
\margin{\uline{Ex.} La corr. \uline{identique} de $\underline E$ sur \unsure{lui-même}.}
Pour cette notion de morphismes, et la \deleted{notion de} composition des correspondances,

\scanpage{62}
\mspage{58}
% archive sheet number: 61

On trouve une catégorie, \deleted{la catégorie} $\unsure{\mathrm{ER}|E\mathcal R|E\mathscr R}$, et ses objets (\unsure{équival\textsuperscript{t}} : des graphes (stricts)), \unsure{mériteraient} un nom, exprimant l'intuition qui s'attache à la situation.
On pourrait appeler un $(E,\unsure{|\mathrm{o}|_E|\lvert\circ\rvert_E})$ un ensemble avec relation de disjonction, un \uline{disjoncteur}, et \deleted{la catégorie \unsure{précéd\textsuperscript{te}}} parler des \uline{morphismes de disjoncteurs}, notation $\unsure{\operatorname{Hom}^{\mathrm{dis}}|\operatorname{Homdis}|\operatorname{Hom}_{\mathrm{dis}}}(\underline E,\underline E')$, et la \uline{catégorie des disjoncteurs}.
Ceci dit, \deleted{\unsure{l'on}} pour un \inserted{tel} morphisme de disjoncteurs
\[
E\xrightarrow{\alpha}\unsure{\mathcal P|\mathscr P|P}(E')
\]
considérons l'application correspondante
\[
\deleted{\unsure{\mathbb\Sigma|\Sigma}_{\unsure{E}}}
\]
\diagram{
\begin{tikzcd}[ampersand replacement=\&,column sep=large]
\overline\alpha:\ \mathcal P(E) \arrow[r] \& \mathcal P(E') \\
A \arrow[r,mapsto] \& \displaystyle\bigcup_{x\in A}\alpha(x)
\end{tikzcd}}
et l'application composée
\diagram{
\begin{tikzcd}[ampersand replacement=\&,column sep=large]
\mathbb\Sigma_{\unsure{\underline E|E}} \arrow[r,hook,"\mathrm{inc}"] \& \mathcal P(E) \arrow[r,"\overline\alpha"] \& \mathcal P(E') \arrow[r,"\operatorname{supp}^{\unsure{\mathrm{o}|\circ|0}}_{E'}"] \& \mathbb\Sigma_{E'}
\end{tikzcd}}
,
\deleted{Il est} on pose
\[
\alpha_*:\mathbb\Sigma_E\longrightarrow\mathbb\Sigma_{E'},\qquad
\alpha_{\unsure{\circ|*}}(S)=\operatorname{supp}^{\mathrm{o}}_{E'}(\overline\alpha(S)).
\]
Cette application est en fait définie \deleted{\unsure{pour}}\inserted{\unsure{quelque|quelle}}

\scanpage{63}
\mspage{59}
% archive sheet number: 62

que soit la correspondance $\alpha$, même si celle-ci n'est pas \uline{compatible} aux données de disjonction. Cependant, quand il y a cette compatibilité, on trouve la formule essentielle
\[
\boxed{\alpha_*(\operatorname{supp}^{\unsure{\mathrm{o}|\circ|0}}_E A)=\operatorname{supp}^{\mathrm{o}}_{E'}(\overline\alpha(A))}\tag*{(*)}
\]
% The following attempted proof is crossed through by long diagonal strokes.
\deleted{\uline{Dém.} Soit $[\overline A=\operatorname{supp}^{\mathrm{o}}_E(A)$, donc $A\subset\overline A]$ alors}
\deleted{$\unsure{\alpha_*(}$}
\[
\deleted{\alpha_*(\overline A)=\operatorname{supp}^{\mathrm{o}}_{E'}(\overline\alpha(\overline A))\supset\operatorname{supp}^{\mathrm{o}}_{E'}(\overline\alpha(A))}
\]
\deleted{puisque $\overline A\supset A$ implique $[\overline\alpha(A)\supset\overline\alpha(\overline A)]$.}
\deleted{Prouvons l'inclusion inverse}
\[
\deleted{\operatorname{supp}^{\mathrm{o}}_{E'}(\overline\alpha(\overline A))\subset\operatorname{supp}^{\mathrm{o}}_{E'}(\overline\alpha(A))},
\]
\deleted{donc que pour tout $x'\in E'$, on a}
\[
\deleted{x'\unsure{|\mathrm{o}||\lvert\circ\rvert}\overline\alpha(A)\Longrightarrow x'|\mathrm{o}|\overline\alpha(\overline A)\quad ?}
\]
\deleted{Je ne vois aucun \unsure{moyen} de le prouver, \unsure{hélas} -- ça doit être faux! Il faudrait une \unsure{correction}. Mais il reste une \uline{partie}.}
On va \unsure{la|le} déduire de la formule ``duale''
\[
\alpha_*(\operatorname{cosupp}^{\mathrm{o}}_E(\deleted{\unsure{A}}\unsure{B|A}))=\operatorname{cosupp}^{\mathrm{o}}_{E'}(\overline\alpha(\deleted{\unsure{A}}\unsure{B|A}))\unsure{\longleftarrow}
\tag*{(**)}
\]
\deleted{\unsure{en effet}}
($\forall B\in\unsure{\mathcal P|\mathscr P|P}(E)$), en appliquant cette dernière \inserted{deux fois \deleted{\unsure{une fois}}} à $B=\operatorname{cosupp}^{\mathrm{o}}_E(A)$\inserted{(une fois pour $B$, une \unsure{autre} fois pour $A$)}.
\[
\begin{aligned}
\alpha_*(\operatorname{supp}^{\mathrm{o}}(\deleted{\unsure{B|A}}\unsure{A|B}))
&=\operatorname{cosupp}^{\mathrm{o}}_{E'}(\overline\alpha(\operatorname{cosupp}^{\mathrm{o}}_E(A)))\\
&\begin{aligned}&\Vert\\
&\operatorname{cosupp}^{\mathrm{o}}_{E'}\bigl(\underbrace{\operatorname{supp}^{\mathrm{o}}_{E'}(\overline\alpha(\text{---}))}_{\substack{\Vert\\\alpha_*(\operatorname{cosupp}^{\mathrm{o}}_E(A))=\deleted{\unsure{\operatorname{cosupp}^{\mathrm{o}}}}}}\bigr)
\end{aligned}
\end{aligned}
\]
% A short isolated dash is visible in the lower left.
-

\scanpage{64}
\mspage{60}
% archive sheet number: 63

et comme $\alpha_*\operatorname{cosupp}^{\unsure{\mathrm{o}|\circ|0}}_E(A)=\operatorname{cosupp}^{\mathrm{o}}_{E'}(\overline\alpha(A))$, on trouve finalement
\[
\begin{aligned}
\operatorname{supp}^{\mathrm{o}}_E(A)&=\operatorname{cosupp}^{\unsure{\mathrm{o}|\circ|0}}_{\unsure{E|E\prime}}(\operatorname{cosupp}^{\mathrm{o}}_{E'}(\overline\alpha(A)))\\
&=\operatorname{supp}^{\mathrm{o}}_{E'}(\overline\alpha(A))
\end{aligned}
\qquad\text{qed.}
\]

Donc il faut établir la formule, qui, pour $B\in\unsure{\mathbb\Sigma|\Sigma}_E$, \deleted{\unsure{donc}} s'écrit
\[
\begin{aligned}
\alpha_*(CB)&=\operatorname{cosupp}^{\mathrm{o}}_{E'}(\overline\alpha(B))\\
&\begin{aligned}&\Vert\\
&\underbrace{\operatorname{cosupp}^{\mathrm{o}}_{E'}\bigl(\underbrace{\operatorname{supp}^{\mathrm{o}}_{E'}(\overline\alpha(B))}_{\substack{\Vert\ \text{déf}\\\alpha_*(B)}}\bigr)}_{C\alpha_*(B)}
\end{aligned}
\end{aligned}
\]
i.e.
\[
\boxed{\alpha_*(CB)=C\alpha_*(B)}
\]
i.e. $\alpha_*$ compatible avec $\unsure{S|C}$.

\deleted{Soit donc} Il faut \unsure{en} prouver $(**)$, qui apparaît comme la formule-clef.
Posons
\[
C=\operatorname{cosupp}^{\mathrm{o}}B=\{x\in E\mid x\unsure{|\mathrm{o}||\lvert\circ\rvert}B\}
\]
alors la formule $(**)$ s'écrit
\[
\operatorname{supp}^{\mathrm{o}}_{E'}(\overline\alpha(C))=\operatorname{cosupp}^{\mathrm{o}}_{E'}(\overline\alpha(B))
\qquad(C=\operatorname{cosupp}^{\mathrm{o}}_E(B)
\]
Prouvons d'abord $\subset$, qui revient à
\[
\overline\alpha(C)\subset\operatorname{cosupp}^{\mathrm{o}}_{E'}(\overline\alpha(B)
\]
ou encore :
\[
\overline\alpha(B)|\mathrm{o}|_{E'}\overline\alpha(\unsure{C|B})
\]
et résulte de \deleted{$B\unsure{\lvert\mathrm{o}|\not\mid}$}

\scanpage{65}
\mspage{61}
% archive sheet number: 64

\[
B\unsure{|\mathrm{o}|_E|\lvert\circ\rvert_E}C\,,
\]
compte tenu de la \deleted{propriété fondamentale}\inserted{compatibilité \unsure{de}} de $\alpha$ aux disjonctions.
Reste la question
\[
\operatorname{cosupp}^{\unsure{\mathrm{o}|\circ|0}}_{E'}(\overline\alpha(B))\subset\operatorname{supp}^{\mathrm{o}}_{E'}(\overline\alpha(C))
\]
i.e. ceci :

$(***)$ Soit $x'\in E'$, alors $x'|\mathrm{o}|\overline\alpha(B)$ implique $x'\in\operatorname{supp}^{\mathrm{o}}_{E'}(\overline\alpha(C))$, i.e. $\forall y'\in E'$, avec $y'|\mathrm{o}|\overline\alpha(C)$, on a $x'|\mathrm{o}|y'$.

Pour le dire autrement :
\[
x',y'\in E',\quad x'|\mathrm{o}|\overline\alpha(B),\ y'|\mathrm{o}|\overline\alpha(C)\Longrightarrow x'|\mathrm{o}|y'
\]

Il ne semble pas que ceci puisse se déduire de la seule compatibilité de $\alpha$ aux disjonctions, qui affirme \unsure{qque|qqque} chose sur la disjonction \inserted{seulement} d'éléments \inserted{$x',y'$} de $E'$ qui \deleted{sont} correspondent à des éléments donnés $x,y$ de $E$.
% slanted
\margin{Non, en effet, voir plus bas}
\deleted{Il semble donc, pour pouvoir conclure}
D'ailleurs, on a un contre-exemple immédiat, en prenant dans $E$ \inserted{et} $E'$ la ``disjonction tautologique'' $x|\mathrm{o}|y\Longleftrightarrow x\ne y$, d'où
$\unsure{\mathbb\Sigma|\Sigma}_E=\unsure{\mathcal P|\mathscr P|P}(E)$, $\mathbb\Sigma_{E'}=\mathcal P(E')$, et une \uline{application} $\alpha:E\longrightarrow E'$, d'où $\alpha_*:\mathbb\Sigma_E\longrightarrow\unsure{\Sigma|\mathbb\Sigma}_{E'}$ donnée par \deleted{$\unsure{\alpha}$} $A\unsure{\longrightarrow|\longmapsto}\alpha(A)$.
\deleted{Même pour cette \unsure{appli}}
La compatibilité de $\alpha$ aux disjonctions signifie ici $\alpha$ injection, mais $\alpha_*$ n'est compatible \unsure{à|:} $C$
